\section{Условие задачи}

\subsection{Цель работы}
Спроектировать иерархию с абстрактным базовым классом, применить динамический полиморфизм и организовать полиморфную коллекцию с ручным управлением памятью.

\subsection{Вариант 19: мир E «Банк»}
Реализовать классы \texttt{Deposit}, \texttt{Credit}, \texttt{Card}. Общая характеристика \texttt{metric()} определяет годовой эффект продукта. Индивидуальное задание --- поиск продукта с максимальной метрикой за один проход.

Обязательны виртуальный деструктор базы, не менее двух чисто виртуальных методов, потоковые операторы, сравнение метрик с точностью $\varepsilon$, приведение к \texttt{double}, специфичные методы наследников и применение \texttt{dynamic\_cast}. Коллекция хранит указатели на базовый класс в собственном динамическом массиве, поддерживает добавление, удаление, обход и суммарную метрику. Используются C++20, CMake 3.20+ и не менее 10 тестов Google Test.
